
\documentclass[12pt]{article}
\usepackage[utf8]{inputenc}
\usepackage[spanish]{babel}
\usepackage{amsmath}
\usepackage{fancyhdr}
\usepackage{geometry} % To set margins
\usepackage{enumitem} % For list customization

% Set page margins
\geometry{a4paper, margin=1in}

% Define custom header
\pagestyle{fancy}
\fancyhf{} % Clear all header and footer fields
\fancyhead[R]{IT1044 (Marrón)\\UTECA 26-1\\Página \thepage{} de 3} % Right header with page number
\renewcommand{\headrulewidth}{0pt} % No horizontal line in header
\renewcommand{\footrulewidth}{0pt} % No horizontal line in footer

% Adjust list spacing to be compact
\setlist[enumerate]{itemsep=0pt,parsep=0pt,topsep=0pt,partopsep=0pt}
\setlist[itemize]{itemsep=0pt,parsep=0pt,topsep=0pt,partopsep=0pt}

\begin{document}

% Content for Page 1
\noindent\textbf{UTECA} \\
\noindent LA UNIVERSIDAD DE/MVS

\vspace{1cm} % Space after the header elements

\noindent\textbf{Título del Curso:} Inglés X \\
\noindent\textbf{Programa y Trayectoria:} Interpretación y Traducción: 10$^\circ$ Cuatrimestre \\
\noindent\textbf{ID del Curso:} IT1044 \\
\noindent\textbf{ID del Cohorte:} 10

\section*{Descripción}
Este curso es la clase final requerida en estudios del idioma inglés para estudiantes del programa de grado en Interpretación y Traducción de la Universidad Tecnológica de América (UTECA). Como tal, es un curso avanzado diseñado para:
\begin{itemize}
    \item perfeccionar y consolidar las técnicas de evaluación presentadas en Inglés IX.
    \item ampliar las bases modernas para la traducción lingüística en inglés.
    \item desarrollar las habilidades prácticas necesarias para la entrada al mundo del empleo profesional.
\end{itemize}
Al igual que en Inglés IX, se utilizará un enfoque constructivista para el aprendizaje avanzado del idioma inglés. Los estudiantes perfeccionarán y consolidarán sus técnicas de evaluación a través del análisis y la crítica de la literatura clásica inglesa y de relatos históricos. Se introducirá a los estudiantes al lenguaje del humor y a los elementos estilísticos de la escritura creativa. Y se introducirá a los estudiantes al lenguaje de la codificación informática mediante aplicación práctica. En preparación para la búsqueda de empleo profesional, los estudiantes evaluarán sitios web de empleo, crearán currículums y cartas de presentación, y practicarán habilidades de entrevista. Aunque la lectura, escritura y el desarrollo de vocabulario seguirán siendo habilidades fundamentales y esenciales, este curso tendrá un mayor enfoque en la presentación. Se revisarán regularmente las estructuras gramaticales.

Las herramientas profesionales de mapeo traslacional desarrolladas en este curso serán específicas para las traducciones del español y el inglés, aunque dichas herramientas y metodologías pueden extenderse fácilmente a la traducción de otros idiomas. Los estudiantes recibirán crédito por cualquier trabajo completado a través de la plataforma de aprendizaje en línea Coursera.

Los temas principales del curso incluirán:
\begin{enumerate}
    \item Literatura y su Crítica
    \item Investigación Académica
    \item La Historia de los Imperios
    \item El Lenguaje de la Codificación y Algunas Aplicaciones Prácticas
    \item Principios de la Escritura Creativa
    \item Elementos y Lenguaje del Humor
    \item Elaboración de Currículums y Cartas de Presentación
    \item Búsqueda de Empleo y Entrevistas
    \item Desarrollo de Vocabulario
\end{enumerate}

\section*{Alcance y Secuencia General}
\noindent I. Literatura y su Crítica
\begin{enumerate}
    \item Repaso de los Elementos de Crítica
    \item Literatura
    \begin{enumerate}
        \item "Joy Luck Club" de Amy Tan
        \item "Brave New World" de Aldous Huxley
    \end{enumerate}
\end{enumerate}

\clearpage % Simulate page break for Page 2

% Content for Page 2
\noindent\textbf{UTECA} \\
\noindent LA UNIVERSIDAD DE/MVS

\vspace{1cm}

\noindent II. Investigación Académica
\begin{enumerate}
    \item Artículos Académicos y su Presentación
    \item Google Scholar
    \item JSTOR
    \item Revistas de Acceso Abierto
\end{enumerate}

\noindent III. La Historia de los Imperios
\begin{enumerate}
    \item El Imperio Han
    \item El Imperio Británico
\end{enumerate}

\noindent IV. El Lenguaje de la Codificación y Algunas Aplicaciones Prácticas
\begin{enumerate}
    \item NetLogo
    \item Python
    \begin{enumerate}
        \item Natural Language Toolkit (NLTK)
        \item APIs
    \end{enumerate}
    \item LaTeX
    \begin{enumerate}
        \item Latexila
        \item Overleaf
    \end{enumerate}
\end{enumerate}

\noindent V. Principios de la Escritura Creativa: Consejos de los Maestros
\begin{enumerate}
    \item Edgar Allan Poe
    \item George Orwell
    \item Kurt Vonnegut
    \item Strunk y White
    \item ... ¡y más!
\end{enumerate}

\noindent VI. Los Elementos y el Lenguaje del Humor
\begin{enumerate}
    \item ¿Por qué es gracioso?
    \begin{enumerate}
        \item Juegos de Lenguaje
        \item Juegos de Realidades
        \item Sátira
    \end{enumerate}
    \item Humor Americano
    \begin{enumerate}
        \item Saturday Night Live (SNL)
        \item Comedias de Situación
    \end{enumerate}
    \item Humor Británico
    \begin{enumerate}
        \item Monty Python
    \end{enumerate}
\end{enumerate}

\noindent VII. Elaboración de Currículums y Cartas de Presentación
\begin{enumerate}
    \item Revisión de Estilos de Currículums
    \item Redacción de un Currículum u Hoja de Vida (curriculum vitae)
    \item Redacción de una Carta de Presentación
\end{enumerate}

\noindent VIII. Búsqueda de Empleo y Entrevistas
\begin{enumerate}
    \item Sitios Web de Empleo
    \begin{enumerate}
        \item Indeed
    \end{enumerate}
\end{enumerate}

\clearpage % Simulate page break for Page 3

% Content for Page 3
\noindent\textbf{UTECA} \\
\noindent LA UNIVERSIDAD DE/MVS

\vspace{1cm}

\noindent VIII. Búsqueda de Empleo y Entrevistas (continuación)
\begin{enumerate}
     \item Sitios Web de Empleo
 %   \setcounter{enumii}{1} % Continue sub-enumeration from previous page
    \begin{enumerate}
        \item Talenteca
        \item LinkedIn
        \item OCCMundial
        \item Portal del Empleo (Bolsa de Trabajo del Gobierno Mexicano)
        \item Bumeran México
    \end{enumerate}
    \item Entrevistas
    \begin{enumerate}
        \item Preparación
        \item Práctica
    \end{enumerate}
\end{enumerate}

\noindent IX. Desarrollo de Vocabulario

\section*{Expectativas}
Se espera que los estudiantes:
\begin{itemize}
    \item Asistan a todas las clases a tiempo
    \item Estén preparados para tomar apuntes y acceder a materiales en línea
    \item Participen en todas las actividades de clase, incluyendo discusiones y presentaciones
    \item Completen todas las tareas y exámenes
\end{itemize}

\section*{Criterios de Egreso}
Tras la finalización exitosa del curso, el estudiante será capaz de:
\begin{itemize}
    \item Analizar, evaluar y proporcionar una crítica escrita de novelas clásicas en inglés
    \item Realizar investigación académica independiente utilizando Google Scholar, JSTOR y revistas de acceso abierto
    \item Realizar investigación histórica independiente, discutir los hallazgos en un trabajo académico y presentar el trabajo verbalmente
    \item Demostrar una comprensión básica de las estructuras de lenguaje de programación y ejecutar algunas aplicaciones prácticas para Python y LaTeX
    \item Escribir una pieza corta de ficción utilizando los principios de la escritura creativa
    \item Definir los elementos básicos del humor americano y británico
    \item Redactar un currículum profesional y una carta de presentación
    \item Navegar por sitios populares de empleo en internet y enviar exitosamente solicitudes de trabajo completadas
    \item Demostrar habilidades de entrevista de trabajo mediante simulación
\end{itemize}

\section*{Evaluaciones Provisionales (Parciales) \hfill Calificación Final}
\noindent\begin{tabular}{ll}
    Trabajo Diario y Participación & 10\% \\
    Tareas & 60\% \\
    Examen Parcial & 30\% \\
\end{tabular}
\hfill
\begin{tabular}{ll}
    1P & 25\% \\
    2P & 25\% \\
    3P & 25\% \\
    Examen Final & 25\% \\
\end{tabular}

\end{document}
